\chapter{What the coupling does}\label{uch:intro}\label{ch:c1}\label{ch:intro}
This chapter says what the coupling does and which pieces take part. It then says what is new for you, whether you
come from Raven or from MODFLOW.

\section{The problem this solves}
Raven \cite{craig2020raven} simulates the water cycle of a watershed: rain and snow, soil moisture, evaporation, runoff,
and flow in rivers and lakes. It divides a watershed into subbasins, and each subbasin into Hydrologic Response Units
(HRUs). An HRU is an area that behaves alike, for example forest on till, or a wetland.

Raven describes groundwater with simple ``buckets'' (linear reservoirs). Water enters a bucket and leaves it at a rate
proportional to what is stored. Buckets are fast and easy to calibrate. But they cannot say \emph{where} groundwater
flows, how deep the water table is at a well, how a pumping well lowers nearby streamflow, or whether a stream gains water
from the aquifer or loses water to it.

MODFLOW~6 \cite{langevin2017mf6} is the U.S. Geological Survey's groundwater model and the international standard for
those questions. It divides the underground into a three-dimensional grid of cells. It calculates the water level (the
\emph{head}) in every cell from the physics of flow through soil and rock. But MODFLOW has no snow, soil or evaporation
processes of its own: it must be told how much water reaches the water table.

The coupling joins the two so that each does what it is good at. Raven calculates everything at and near the land
surface. MODFLOW calculates the saturated groundwater below. The two exchange water every day.

\begin{plain}
Raven decides how much water soaks down to the water table. MODFLOW moves that water sideways and downwards through
the aquifer. Wherever the aquifer overflows --- into rivers, wetlands, lakes or back to the soil --- MODFLOW hands the
water back to Raven. Pumping wells and other stresses act on MODFLOW's side. Both models keep exact accounts, so not
a drop is created or lost at the boundary between them.
\end{plain}

\section{Two ways to use the coupling}\label{sec:twoways}
Most users are in one of two situations (Figure~\ref{fig:twoways}).
\begin{itemize}
\item \textbf{You already have a MODFLOW~6 model}, perhaps one calibrated over years. Raven connects to it and uses it as
      it is: same grid, same layers, same wells and boundaries (Chapter~\ref{uch:existing}).
\item \textbf{You have a Raven model but no groundwater model.} Raven builds one for you from your HRUs and a short
      description of the soil and rock layers under them (Chapter~\ref{uch:build}).
\end{itemize}
Everything else is the same in both cases: how water is exchanged, how the accounts are kept, the outputs and the tests.

\begin{figure}[!htbp]\centering% Which chapter to start from (TikZ; text is set by LaTeX, so boxes always fit their text)
\begin{tikzpicture}[font=\sffamily\small,
  base/.style={rounded corners=3pt,line width=0.8pt,align=center,inner sep=5pt,execute at begin node={\hyphenpenalty=10000\exhyphenpenalty=10000}},
  q/.style={base,draw=brand,fill=white,text width=58mm,font=\sffamily\bfseries},
  yesA/.style={base,draw=teal!70!black,fill=teal!7,text width=66mm},
  noA/.style={base,draw=orange!80!black,fill=orange!7,text width=66mm},
  yesB/.style={base,draw=teal!70!black,fill=white,text width=66mm,font=\sffamily\footnotesize},
  noB/.style={base,draw=orange!80!black,fill=white,text width=66mm,font=\sffamily\footnotesize},
  fin/.style={base,draw=brand,fill=brandlight,text width=144mm},
  arr/.style={-{Stealth[length=2.4mm]},line width=0.9pt},
  lab/.style={font=\sffamily\bfseries\small,inner sep=2pt}]
\node[q] (q) at (0,0) {Do you already have a MODFLOW~6 groundwater model?};
\node[yesA,anchor=north] (y) at (-39mm,-11mm) {\textbf{Raven connects to your model}\\(Chapter~\ref{uch:existing})\\[3pt]Your model is used as it is. Raven links its HRUs to the model's top layer.};
\node[noA,anchor=north] (n) at (39mm,-11mm) {\textbf{Raven builds a groundwater model}\\(Chapter~\ref{uch:build})\\[3pt]From your HRUs and a short description of the soil and rock layers under them.};
\node[yesB,anchor=north] (y2) at ([yshift=-5mm]y.south) {Tell Raven where the model is on the map: its projection (\texttt{:MF6CRS EPSG:\dots}), a file of cell outlines, or a table of overlap weights.};
\node[noB,anchor=north] (n2) at ([yshift=-5mm]n.south) {Choose a grid, or keep the default: regular, rotated, quadtree, cells that follow HRU boundaries, or your own cells.};
\path let \p1=(y2.south), \p2=(n2.south) in coordinate (m) at (0,{min(\y1,\y2)-7mm});
\node[fin,anchor=north] (e) at (m) {Then, every day of the run, Raven sends recharge down and MODFLOW sends seepage and river water back up. Both keep exact water accounts (Chapters~\ref{uch:exchanges} and~\ref{uch:outputs}).};
\draw[arr,teal!70!black] (q.west) -| node[lab,pos=0.35,above,text=teal!70!black]{Yes} (y.north);
\draw[arr,orange!80!black] (q.east) -| node[lab,pos=0.35,above,text=orange!80!black]{No} (n.north);
\draw[arr,teal!70!black] (y.south) -- (y2.north);
\draw[arr,orange!80!black] (n.south) -- (n2.north);
\draw[arr,teal!70!black] (y2.south) -- (y2.south |- e.north);
\draw[arr,orange!80!black] (n2.south) -- (n2.south |- e.north);
\end{tikzpicture}

\caption{Which chapter to start from: connecting Raven to an existing MODFLOW~6 model, or letting Raven build one.}\label{fig:twoways}
\end{figure}

\section{How to use the two volumes}
The manual has two volumes. This one, the \emph{User Manual}, tells you what to write and what to expect. Volume~2,
\emph{Reference and Verification} (\cmd{Raven\_MF6\_Reference.pdf}), holds every rule and equation, the full tables of
commands and messages, the tests and the developer guide. In this volume, ``Reference, Section~4.3'' means Section~4.3
of Volume~2. Table~\ref{tab:readguide} says where to go.
\begin{table}[!htbp]\centering\small
\caption{Where to find what.}\label{tab:readguide}
\begin{tabularx}{\linewidth}{P{0.34\linewidth}Y}\toprule\hdr
If you want to\dots & Read\\\midrule
understand the idea & this chapter; groundwater in plain terms: Section~\ref{sec:newraven} and Reference, Chapter~\ref{R-ch:c2}\\
run a first coupled model & Chapter~\ref{uch:quick} (example \cmd{Liard\_groundwater})\\
let Raven build a groundwater model & Chapter~\ref{uch:build}; example \cmd{Liard\_groundwater}\\
connect Raven to your MODFLOW model & Chapter~\ref{uch:existing}; example \cmd{Liard\_existing\_model}\\
know what water moves where & Chapter~\ref{uch:exchanges}; the equations: Reference, Chapters~\ref{R-ch:c4} and~\ref{R-ch:c5}\\
read the results and check the water balance & Chapter~\ref{uch:outputs}\\
restart, run ensembles or calibrate & Chapter~\ref{uch:restarts}\\
fix a problem & Chapter~\ref{uch:trouble}; every message: Reference, Sections~\ref{R-app:errors} and~\ref{R-app:warnings}\\
look up a command or an output column & Chapter~\ref{uch:commands}; Reference, Chapter~\ref{R-ch:c6} and Appendix~\ref{R-app:tables}\\
see how the code was tested & Reference, Chapter~\ref{R-ch:c8}\\
change the code & Reference, Chapter~\ref{R-ch:c9}\\\bottomrule
\end{tabularx}
\end{table}
Technical words are explained in the glossary (Appendix~\ref{app:uglossary}). The workshops in
\cmd{docs/tutorials/workshops} of RavenHydroFramework-extra take you through the same steps with your own hands: Workshop~1
(getting started), 2 (your own groundwater model), 3 (the soil zone), 4 (an existing MODFLOW~6 model), 5 (restarts,
ensembles and particle tracking) and 6 (calibration with Ostrich).

\section{What you do, and what Raven does}
You work only with Raven's familiar text files:
\begin{enumerate}
\item mark which HRUs have groundwater underneath (\cmd{AQUIFER\_PROFILE} in the \cmd{.rvh} file);
\item describe the aquifer materials and their layering (\cmd{.rvp} file);
\item give the map of HRUs (a GeoJSON file) and, if you have them, land-surface and bedrock maps, rivers, wells,
monitoring wells and boundaries (the \cmd{.rvg} file, which the coupling adds to Raven's inputs);
\item switch the coupling on (\cmd{:GroundwaterModel MODFLOW6} in the \cmd{.rvi} file).
\end{enumerate}
Raven then does everything else by itself. It lays a grid over the HRUs and works out how much of each HRU lies in each
grid cell. It builds the layers from the land surface down to bedrock and writes a standard MODFLOW~6 model. It starts
MODFLOW, exchanges water every time step, and writes the results next to its usual outputs.

MODFLOW is used as a \emph{library}: a program part that another program can call. Raven talks to it through MODFLOW's
official programming interface \cite{hughes2022xmi}. No MODFLOW source code is changed.

\section{What is new for a Raven user}\label{sec:newraven}
If you know Raven but not MODFLOW, these are the new ideas. Reference, Chapter~\ref{R-ch:c2}, explains them with
equations and numbers.
\begin{itemize}
\item \textbf{Head.} The level to which water rises in a well, in metres above sea level. Groundwater flows from high
head to low head. Near the surface the head is the water table.
\item \textbf{Cells and layers.} MODFLOW computes one head per cell. The cells form layers, from the model top down to
bedrock. A column is the stack of cells at one map location.
\item \textbf{Hydraulic conductivity} ($K$, m/d) says how easily water moves through a material: gravel passes water
fast, clay very slowly. Flow between two cells is the head difference times a conductance, which comes from $K$, the
cell sizes and the thickness of saturated material.
\item \textbf{Storage.} When the water table falls, the pores drain; the specific yield ($S_y$) is the share of the volume
that drains. A full (confined) layer releases far less water, through the specific storage ($S_s$).
\item \textbf{Boundaries.} Rivers, drains, wells and fixed heads add or remove water at chosen cells. Many of them depend
on the head: a river leaks into the aquifer when the head is below the river, and gains water when it is above.
\item \textbf{Your \cmd{GROUNDWATER} store changes role.} In a coupled HRU it is a pass-through: whatever enters it during
a day is handed to MODFLOW as recharge that day. Baseflow then comes from the aquifer, not from a bucket.
\end{itemize}

\section{What is new for a MODFLOW user}\label{sec:newmodflow}
If you know MODFLOW but not Raven, these are the new ideas.
\begin{itemize}
\item \textbf{HRUs and subbasins.} Raven's land units are HRUs, grouped into subbasins; each subbasin drains to one river
reach. Raven computes each HRU's water budget, then routes the runoff down the river network.
\item \textbf{State variables.} Raven keeps water in named stores per HRU: \cmd{SOIL[0]}, \cmd{SOIL[1]}, \dots,
\cmd{DEPRESSION}, \cmd{SURFACE\_WATER}, \cmd{GROUNDWATER} and others. Processes in the \cmd{.rvi} file move water between
them; for example \cmd{:Percolation PERC\_CONSTANT FAST\_RESERVOIR GROUNDWATER} moves water from one store to
another.
\item \textbf{Recharge is computed, not given.} Raven's snow, soil and evaporation processes decide how much water
reaches the water table each day. Your RCH and EVT packages are replaced by Raven's (Chapter~\ref{uch:existing}).
\item \textbf{The files.} \cmd{.rvi} holds the options and processes, \cmd{.rvh} the subbasins and HRUs, \cmd{.rvp} the
parameters, \cmd{.rvt} the time series (forcing, well rates, observations), \cmd{.rvc} the initial conditions. The
coupling adds the \cmd{.rvg} file (Table~\ref{tab:ufiles}).
\item \textbf{Time.} Raven steps through time, usually one day at a time. MODFLOW runs one Raven step per MODFLOW time
step; your stress periods are kept and split into Raven steps.
\item \textbf{Calibration.} Raven is often calibrated with Ostrich against streamflow; with the coupling, heads at
monitoring wells join the objective.
\end{itemize}

\section{The pieces}\label{ch:overview}
\subsection{Who owns which water}
Every drop of water in the coupled system belongs to exactly one of the two models at any time:
\begin{itemize}
\item \textbf{Raven} owns snow, canopy, soil and vadose-zone stores, depressions and wetlands, rivers, lakes and
reservoirs.
\item \textbf{MODFLOW} owns the saturated groundwater below the \emph{model top}.
\end{itemize}
By default the model top is the land surface. Users with physically based soil layers can lower it with
\cmd{SOIL\_ZONE\_DEPTH}. Or, with \cmd{:SoilZone PHYSICAL}, they can hand the soil layers to Raven and let the water table
rise into them, as in GSFLOW \cite{markstrom2008gsflow} (Section~\ref{sec:usoilzone}).

On sloping and irregular land each cell then samples the DEM at many points. So the soil saturates first at the foot of
a slope and in hollows, and the saturated area spreads up and retreats down the slope with the water table
(Section~\ref{sec:applysz}; the rules: Reference, Section~\ref{R-sec:szslope}). Below the soil, recharge can cross an
unsaturated zone (Richards' equation, a kinematic wave or a storage cascade), and dry soil draws water up from a shallow
water table by capillary rise (Reference, Sections~\ref{R-sec:szvadose} and~\ref{R-sec:szcaprise}).

Water moves between the two models only through the exchanges of Chapter~\ref{uch:exchanges}, and each exchange is
recorded on both sides. In a coupled HRU, Raven's \cmd{GROUNDWATER} store is a \emph{pass-through}: whatever enters it
during a day is handed to MODFLOW as recharge that day.

\subsection{The software}
Figure~\ref{fig:arch} shows the parts. The \emph{model builder} runs once at the start: it reads the inputs and writes a
standard MODFLOW~6 model. The \emph{exchange} runs every time step.
\begin{figure}[!htbp]\centering
\begin{tikzpicture}[box/.style={draw,rounded corners,align=center,minimum height=11mm,text width=34mm,font=\small},
 ar/.style={-{Stealth[length=2.2mm]},thick}]
\node[box,fill=teal!12] (in) at (0,1.1) {Raven inputs\\ \cmd{rvi rvh rvp rvg rvc rvt}};
\node[box,fill=gray!12] (geo) at (0,-1.1) {Maps: HRU polygons, DEM, bedrock, rivers};
\node[box,fill=gray!20] (bld) at (5,0) {Model builder\\ grid, layers, links,\\ MODFLOW 6 files};
\node[box,fill=violet!12] (mf) at (10,0) {\cmd{libmf6}\\ MODFLOW 6};
\node[box,fill=teal!12] (rv) at (10,-3.2) {Raven solver\\ HRU processes, routing};
\node[box,fill=teal!6] (out) at (5,-3.2) {Outputs\\ diagnostics, hotstart,\\ budget, heads};
\draw[ar] (in) -- (bld); \draw[ar] (geo) -- (bld);
\draw[ar] (bld) -- node[above,font=\scriptsize]{once} (mf);
\draw[ar,<->] (mf) -- node[right,font=\scriptsize,align=left]{water and heads\\ every time step} (rv);
\draw[ar] (rv) -- (out);
\end{tikzpicture}
\caption{The software pieces. The builder runs once per simulation (or per ensemble member); the exchange runs every
time step.}\label{fig:arch}
\end{figure}

\subsection{One day of the coupled model}
This happens every time step (one day in the examples). Reference, Section~\ref{R-ch:timestep}, gives the exact
algorithm.
\begin{enumerate}
\item \textbf{Raven's land-surface day.} Raven calculates snowmelt, infiltration, evaporation and flow between its soil
layers. In coupled HRUs, water that percolates past the soil ends up in \cmd{GROUNDWATER}.
\item \textbf{Recharge handed over.} The water in each HRU's \cmd{GROUNDWATER} is spread over the grid cells under that
HRU, in proportion to how much of the HRU lies in each cell.
\item \textbf{Stresses set.} River levels (from Raven's rivers at the start of the day), well pumping rates, boundary
heads, wetland and lake levels, and the evaporation demand not met by Raven are passed to MODFLOW.
\item \textbf{MODFLOW's day.} MODFLOW calculates new heads everywhere, with the flows to and from rivers, drains,
wells and boundaries.
\item \textbf{Water handed back.} Groundwater that seeps out at the land surface is added to a Raven store of the HRUs
above. Groundwater that discharges to rivers is added to the river inflow of the subbasin. River losses are taken from
the river. Wetland and lake exchanges change those Raven stores.
\item \textbf{Raven's river day.} Raven routes water down the river network, with the groundwater contributions.
Reservoirs compute their seepage using the aquifer head beneath them.
\item \textbf{Accounts.} Both models update their water balances, and the groundwater outputs are written.
\end{enumerate}

\section{Main features}
\begin{itemize}
\item Any number of \emph{aquifer profiles} (layer stacks), each with any number of layers. Layer types: aquifer,
confined aquifer, aquitard (confining layer) and aquiclude.
\item Land surface and bedrock from raster maps. Layers that thin out against bedrock are handled automatically.
\item Coupling to an existing MODFLOW~6 model. Raven links its HRUs to the model's top active layer and leaves the model
intact, with built-in map projections, unit conversion and safeguards against double counting
(Chapter~\ref{uch:existing}).
\item Five grid options: the default regular grid, rotated and refined regular grids, quadtree grids, meshes whose cells
follow HRU boundaries, and imported cells. Layers can be connected by elevation (Section~\ref{sec:ugrids}).
\item Water exchanges: recharge (with optional delay), groundwater seepage at the land surface, rivers (gaining and
losing), evaporation from a shallow water table, wetlands and lakes held in Raven stores, Raven reservoirs, pumping and
injection wells, head-dependent and fixed-head boundaries.
\item Cells can dry out and rewet; areas of aquifer can disconnect and reconnect by themselves.
\item Monitoring wells: simulated heads are written out and scored against observations like streamflow, so the
groundwater parameters can be calibrated (for example with Ostrich).
\item Hotstart (continuing a run from where another ended), ensembles, NetCDF output of the whole head field.
\item Exact water accounting on both sides, checked automatically by an audit tool supplied with the code.
\end{itemize}

\section{Design principles}
\begin{enumerate}
\item \textbf{Raven is the single user interface.} All inputs are Raven files; MODFLOW files are generated.
\item \textbf{Standard MODFLOW only.} Only standard MODFLOW~6 packages are used; the generated model can be opened in
FloPy or ModelMuse.
\item \textbf{Each fact is entered once.} Geometry comes from maps, hydrogeology from profiles, time-varying stresses
from time series.
\item \textbf{Generic.} Nothing is specific to a basin; any Raven model can couple any subset of its HRUs.
\item \textbf{Exact accounting.} Every exchanged volume is computed once and applied identically on both sides.
\item \textbf{No surprise for existing users.} Models without groundwater give bit-for-bit the same results as
Raven~4.15 (the version this coupling is built on), except in the few cases where the coupled code corrects Raven
itself: models with reservoirs, ensembles and age-tracer runs started from a hotstart file (Reference,
Table~\ref{R-tab:ravendiff}). Raven still builds and runs without MODFLOW installed.
\end{enumerate}

\section{Scope and assumptions}
\begin{itemize}
\item \textbf{What MODFLOW simulates:} saturated groundwater flow (Darcy flow, storage change) in a layered aquifer
system under the coupled HRUs, with dry and rewetting cells.
\item \textbf{What stays in Raven:} everything above the model top --- snow, canopy, soil and unsaturated zone, surface
runoff, depressions, rivers, lakes and reservoirs --- and all HRUs without an aquifer profile.
\item \textbf{Water density and temperature} are constant. There is no salt-water or heat transport.
\item \textbf{No constituent transport.} The coupling exchanges water, not constituent mass. So Raven's \cmd{:Transport}
(tracers, nutrients, the age tracer, temperature) is not available in a coupled run; the run stops with a message.
\item \textbf{The grid} is regular (square cells, optionally rotated), a quadtree, an HRU mesh whose cells follow HRU
boundaries, or imported cell polygons (Section~\ref{sec:ugrids}). An existing MODFLOW~6 model keeps its own DIS or DISV
grid (Chapter~\ref{uch:existing}). A generated model covers the cells that the coupled HRUs cover.
\item \textbf{Time:} both models use Raven's time step (daily or shorter). Exchanges are explicit between the models at
the step level, but every head-dependent exchange is implicit inside MODFLOW. With \cmd{:SoilZone PHYSICAL} the
soil--aquifer interface is also iterated inside MODFLOW's outer iterations (Reference, Section~\ref{R-sec:szinterface}).
\end{itemize}

\section{When to use the coupling, and when not to}
Use it when the question involves the \emph{spatial} behaviour of groundwater: the water-table depth at a location, the
effects of pumping on streams, gaining and losing reaches, groundwater flow between subbasins, wetlands fed by
groundwater, the drought recession of large aquifers, or when head measurements should constrain the model.

Raven's conceptual baseflow remains the better choice for fast screening, very large ensembles, or basins without a
significant aquifer.

\section{What changes for an existing Raven model}
\begin{itemize}
\item Nothing, unless \cmd{:GroundwaterModel MODFLOW6} is added. Results then stay bit-for-bit identical to Raven~4.15,
except where the coupled code corrects Raven itself (Reference, Table~\ref{R-tab:ravendiff}):
  \begin{itemize}
  \item models with reservoirs: absolute-stage reservoirs without an initial stage start at their crest, a drying
  reservoir books only the water it has, and \cmd{solution.rvc} carries stages and flows with 12 digits;
  \item ensembles: every member starts from the first run's state;
  \item age-tracer runs started from a hotstart file: the ages continue from the run that wrote the file.
  \end{itemize}
\item With the coupling on, coupled HRUs lose their conceptual deep-groundwater behaviour: their \cmd{GROUNDWATER} store
becomes a pass-through, and their baseflow comes from the aquifer via seepage and river exchange.
\item Run times increase. The 20-year Liard run (3\,252 active columns) takes about 10 minutes instead of 35 seconds on
one processor core.
\end{itemize}

\section{Files at a glance}
\begin{table}[!htbp]\centering
\caption{Input and output files of the coupling.}\label{tab:ufiles}
\begin{tabular}{P{0.2\linewidth}P{0.748\linewidth}}\toprule\hdr
File & Role in the coupling\\\midrule
\cmd{.rvi} & switch on (\cmd{:GroundwaterModel MODFLOW6}); route recharge into \cmd{GROUNDWATER}\\
\cmd{.rvh} & \cmd{AQUIFER\_PROFILE} per HRU; reservoirs with \cmd{:SeepageParameters}\\
\cmd{.rvp} & aquifer classes, profiles and profile parameters\\
\cmd{.rvg} & geometry (HRU polygons, maps, rivers), grid, wells, monitoring wells, boundaries, options\\
\cmd{.rvt} & well rates, boundary heads, observed heads\\
\cmd{.rvc} & initial heads; hotstart blocks written by Raven\\
GeoJSON, ASCII grids & HRU polygons, river lines, boundary shapes; DEM and bedrock maps\\
\cmd{<output>/mf6/} & generated MODFLOW 6 model and its own outputs\\\bottomrule
\end{tabular}
\end{table}
